\documentclass[a4paper]{article}

\usepackage[a4paper, top=8mm, bottom=8mm, left=8mm, right=8mm]{geometry}

\usepackage{fontspec}
\setmainfont{CMU Serif}[Ligatures=TeX] 
\newfontfamily\cyrillicfont{CMU Serif}[Ligatures=TeX]
\newfontfamily{\russianfonttt}[Scale=0.7]{DejaVuSansMono}

\usepackage{polyglossia}
\setmainlanguage{russian}
\setotherlanguage{english}

\usepackage[tiny, compact]{titlesec}

\usepackage{titling}
\setlength{\droptitle}{-1cm}
\pretitle{\begin{center}\begin{bfseries}\Large}
\posttitle{\par\end{bfseries}\end{center}}
\preauthor{\begin{center}\normalsize}
\postauthor{\par\end{center}\vspace{-1.8cm}}

\usepackage{amsmath}
\usepackage{csquotes}
\usepackage{booktabs}
\usepackage{multirow}

\usepackage{hyperref}
\usepackage{bookmark}

\title{Реализация алгоритмов подсчета треугольников в графе на основе GraphBLAS API}

\author{Сергеев Дмитрий Сергеевич}

\date{}

\begin{document}
\maketitle
\begin{flushright}
\begin{tabular}{r}
    Группа: \emph{\enquote{25.Б41-мм}} \\
    Кафедра: \emph{системного программирования} \\
    Научный руководитель: \emph{Григорьев Семен Вячеславович} \\
    Номер семестра практики: \emph{2} \\
\end{tabular}
\end{flushright}

\section{Постановка задачи}

Подсчёт треугольников в графе~--- одна из базовых операций анализа графов, используемая, например, для вычисления кластерного коэффициента и анализа социальных сетей.
Представление графовых алгоритмов в терминах линейной алгебры над разреженными матрицами позволяет переиспользовать высокооптимизированные вычислительные движки, что и реализуется в стандарте GraphBLAS.

Работа выполнялась в рамках учебной ознакомительной практики и направлена на знакомство с GraphBLAS API и сравнительный анализ производительности существующих библиотечных решений на языках Python и C.

\textbf{Целью работы является} реализация алгоритмов подсчета треугольников в графе на основе GraphBLAS API, экспериментальное сравнение собственных реализаций между собой, а также с готовыми реализациями из библиотеки LAGraph.

Для достижения поставленной цели были решены следующие задачи:
\begin{itemize}
    \item выполнить обзор подходов к интеграции GraphBLAS API на языках C и Python;
    \item реализовать алгоритмы подсчёта треугольников с использованием Python-GraphBLAS;
    \item провести замеры производительности разработанных решений и готовых реализаций из LAGraph;
    \item локализовать причины различий в производительности с помощью профилирования.
\end{itemize}

\section{Описание предлагаемого решения}

В качестве вычислительного ядра как в Python-версии (\texttt{python-graphblas}), так и в C-версии (\texttt{LAGraph}), используется библиотека SuiteSparse:GraphBLAS версии 9.4.5, что делает сравнение корректным~--- обе реализации опираются на одну и ту же программную основу.

Были реализованы три алгоритма подсчёта треугольников с использованием \texttt{python-graphblas}:
\begin{itemize}
    \item \textbf{Алгоритм Burkhard} основан на вычислении количества треугольников по формуле:
    \[
    \text{triangles} = \frac{1}{6} \sum_{i} \sum_{j} \left( (A^2)_{ij} \cdot A_{ij} \right)
    \]
    где $A$~--- матрица смежности. В терминах GraphBLAS данная операция представляет собой перемножение матриц $A \times A$ с использованием маски $A$.
    
    \item \textbf{Алгоритм Sandia} работает с нижнетреугольной частью матрицы смежности и использует формулу:
    \[
    \text{triangles} = \sum_{i} \sum_{j} \left( (L \times L)_{ij} \cdot L_{ij} \right)
    \]
    где $L = \text{tril}(A)$.
    
    \item \textbf{Наивный алгоритм} вычисляет количество треугольников через след куба матрицы смежности:
    \[
    \text{triangles} = \frac{\text{tr}(A^3)}{6}
    \]
    где $\text{tr}(A) = \sum_{i} a_{ii}$~--- след матрицы (trace), представляющий собой сумму всех элементов, стоящих на её главной диагонали.
\end{itemize}
\section{Экспериментальное исследование производительности}

Замеры проводились на ноутбуке Huawei MCLF-XX с процессором Intel Core i5-12450H и 16~ГБ оперативной памяти; версия gcc~13.3, Python~3.12.3, ОС Ubuntu~24.04.4.
Для стабилизации окружения были приняты следующие меры: отключены фоновые системные сервисы и таймеры, отключён режим Turbo Boost, отключены состояния простоя CPU (C-states), обработка прерываний перенесена на E-cores, замеры привязаны к P-cores через \texttt{taskset -c 0-7}.
Замеры производились с питанием от сети, порог заряда был ограничен диапазоном 85--90\%.
Для каждого алгоритма на каждом графе проводилось 45 запусков, из которых первые 15 считались прогревочными и не учитывались в результатах (для \texttt{python-graphblas}).
Перед каждой серией замеров выполнялся сброс файлового кеша ОС.

В качестве входных данных использовались графы из коллекции SuiteSparse Matrix Collection различных размеров и структур~--- от плотных графов покупок до разреженных дорожных сетей, чтобы оценить устойчивость результатов на разной топологии (таблица~\ref{tab:datasets}).

\begin{table}[h!]
    \centering
    \caption{Используемые наборы данных.}
    \label{tab:datasets}
    \begin{tabular}{lrr}
        \toprule
        \textbf{Граф} & \textbf{Вершины} & \textbf{Рёбра} \\
        \midrule
        amazon0302    & 262\,111    & 1\,234\,877  \\
        amazon-2008   & 735\,323    & 5\,158\,388  \\
        cit-Patents   & 3\,774\,768 & 16\,518\,948 \\
        roadNet-CA    & 1\,971\,281 & 2\,766\,607  \\
        web-NotreDame & 325\,729    & 1\,497\,134  \\
        web-Stanford  & 281\,903    & 2\,312\,497  \\
        \bottomrule
    \end{tabular}
\end{table}

Наивный алгоритм был исключён из основных замеров: на использованных графах процесс завершался с ошибкой нехватки памяти.

Для каждой комбинации (граф, алгоритм, реализация) проводилась статистическая обработка результатов: тесты на нормальность, среднее, стандартное отклонение, стандартная ошибка среднего (SEM) и доверительный интервал с уровнем доверия 95\%.
Результаты приведены в таблице~\ref{tab:benchmarks}.

\begin{table}[h!]
    \centering
    \small
    \caption{Время выполнения алгоритмов подсчёта треугольников, в секундах, Confidence Interval = 0.95.}
    \label{tab:benchmarks}
    \begin{tabular}{llcc}
        \toprule
        \textbf{Граф} & \textbf{Конфигурация} & \textbf{Время} & \textbf{Std (\%)} \\
        \midrule
        \multirow{4}{*}{web-NotreDame}
            & burkhard-py   & $0.126 \pm 0.001$ & 1.57 \\
            & sandia-py     & $0.031 \pm 0.001$ & 4.85 \\
            & burkhard-lagr & $0.124 \pm 0.001$ & 2.16 \\
            & sandia-lagr   & $0.033 \pm 0.001$ & 2.64 \\
        \midrule
        \multirow{4}{*}{web-Stanford}
            & burkhard-py   & $1.178 \pm 0.001$ & 0.52 \\
            & sandia-py     & $0.264 \pm 0.001$ & 0.51 \\
            & burkhard-lagr & $1.180 \pm 0.004$ & 0.87 \\
            & sandia-lagr   & $0.261 \pm 0.001$ & 1.00 \\
        \midrule
        \multirow{4}{*}{roadNet-CA}
            & burkhard-py   & $0.058 \pm 0.001$ & 0.57 \\
            & sandia-py     & $0.046 \pm 0.001$ & 1.30 \\
            & burkhard-lagr & $0.056 \pm 0.001$ & 0.99 \\
            & sandia-lagr   & $0.054 \pm 0.001$ & 1.32 \\
        \midrule
        \multirow{4}{*}{cit-Patents}
            & burkhard-py   & $4.322 \pm 0.012$ & 0.95 \\
            & sandia-py     & $0.679 \pm 0.001$ & 0.30 \\
            & burkhard-lagr & $4.309 \pm 0.005$ & 0.29 \\
            & sandia-lagr   & $0.698 \pm 0.001$ & 0.42 \\
        \midrule
        \multirow{4}{*}{amazon-2008}
            & burkhard-py   & $0.294 \pm 0.001$ & 0.44 \\
            & sandia-py     & $0.103 \pm 0.001$ & 0.29 \\
            & burkhard-lagr & $0.302 \pm 0.001$ & 0.29 \\
            & sandia-lagr   & $0.111 \pm 0.001$ & 0.86 \\
        \midrule
        \multirow{4}{*}{amazon0302}
            & burkhard-py   & $0.087 \pm 0.001$ & 0.38 \\
            & sandia-py     & $0.027 \pm 0.001$ & 0.87 \\
            & burkhard-lagr & $0.084 \pm 0.001$ & 0.65 \\
            & sandia-lagr   & $0.025 \pm 0.001$ & 3.57 \\
        \bottomrule
    \end{tabular}
\end{table}

Выбросы наблюдались в части серий замеров, тест на нормальность не проходился.
Однако выбросы присутствуют у обеих реализаций примерно в равной мере и не меняют относительного положения доверительных интервалов, поэтому на выводы они не влияют.

Алгоритм Sandia быстрее алгоритма Burkhard на всех графах из таблицы~\ref{tab:benchmarks}: разница составляет от 2 до 6~раз, доверительные интервалы не пересекаются.
Это объясняется тем, что Sandia работает с нижнетреугольной матрицей $L$, содержащей примерно вдвое меньше ненулевых элементов, чем полная матрица $A$.

Реализация на \texttt{python-graphblas} не уступает реализации LAGraph по времени работы.
На большинстве графов конфигурации \texttt{burkhard-py}/\allowbreak\texttt{burkhard-lagr} и \texttt{sandia-py}/\allowbreak\texttt{sandia-lagr} показывают статистически неразличимое время~--- доверительные интервалы пе\-ре\-се\-ка\-ют\-ся (см. таблицу~\ref{tab:benchmarks}, например web-NotreDame, web-Stanford, cit-Patents, amazon-2008).
На грахах roadNet-CA и amazon0302 реализация на Python оказалась быстрее LAGraph для алгоритма Sandia.

Для выяснения причины превосходства \texttt{python-graphblas} на графах roadNet-CA, amazon-2008 и web-NotreDame было проведено профилирование с помощью \texttt{perf} (алгоритм Sandia, частота дискретизации 99~Гц). Вклад внутренних функций GraphBLAS рас\-счи\-ты\-вал\-ся как отношение числа зарегистрированных отсчетов конкретной функции к общему числу отсчетов в точке входа в измеряемый цикл вычислений алгоритма.

Профилирование показало следующее.
\begin{itemize}
    \item На графе roadNet-CA разница наиболее выражена: функция \texttt{GB\_AxB\_saxpy3\_flopcount} занимает $\approx 19{,}08\%$ времени у LAGraph против $\approx 10{,}41\%$ у \texttt{python-graphblas}; \texttt{GB\_select\_positional\_phase1}~--- $\approx 14{,}96\%$ против $\approx 7{,}07\%$; \texttt{GB\_msort\_3}~--- $\approx 6{,}59\%$ против $\approx 3{,}29\%$.
    \item На графе amazon-2008 разница между реализациями минимальна (в пределах 1\% на каждую функцию), однако функции, вызываемые этими операциями, также занимают чуть больше времени у LAGraph, что в сумме даёт суммарную разницу.
    \item На графе web-NotreDame: \texttt{GB\_AxB\_saxpy3\_flopcount}~--- $\approx 4{,}83\%$ против $\approx 2{,}74\%$, \texttt{GB\_msort\_3}~--- $\approx 5{,}22\%$ против $\approx 2{,}74\%$.
\end{itemize}

Во всех трёх случаях LAGraph стабильно тратит больше относительного времени на функции \texttt{GB\_AxB\_sax\-py3\_flop\-count}, \texttt{GB\_select\_posi\-ti\-o\-nal\_pha\-se1} и \texttt{GB\_msort\_3}.
Точная причина различия (например, отличия во внутреннем формате матрицы, передаваемой из Python и из C) не уста\-нав\-ли\-ва\-лась~--- зафиксирован лишь устойчивый факт различия в относительном времени выполнения одних и тех же внутренних функций SuiteSparse.

Полученные результаты подтверждают, что реализация на \texttt{python-graphblas} может не уступать готовым реализациям из LAGraph, а на части графов даже превосходить их.

\section{Заключение}

В ходе выполнения работы получены следующие результаты:

\begin{itemize}
    \item выполнен обзор подходов к интеграции GraphBLAS API на языках C и Python, выбрана библиотека SuiteSparse:GraphBLAS~9.4.5 как общий вычислительный движок для сравнения;
    \item реализованы алгоритмы подсчёта треугольников (Burkhard, Sandia, наивный) с использованием Python-GraphBLAS API;
    \item проведены замеры производительности разработанных решений и готовых реализаций из LAGraph: Sandia быстрее Burkhard в 2--6~раз на всех наборах данных, реализация на \texttt{python-graphblas} не уступает LAGraph, а на части графов превосходит её;
    \item с помощью профилирования (\texttt{perf}, flame graphs) локализованы функции SuiteSparse, на которые расходуется больше времени в LAGraph по сравнению с \texttt{python-graphblas}.
\end{itemize}

Репозиторий: \url{https://github.com/fox56tm/graph_blas_triag_counting}.

Скрипты для подготовки окружения и запуска замеров, инструкции по сборке и воспроизведению результатов, а также построенные боксплоты и flame graphs представлены в файле \texttt{README.md} репозитория.

\end{document}
